\section{Related Work}

Diffusion generative models were first proposed by \citet{sohl2015deep} and have gained renewed attention recently due to strong results on image and waveform generation \citep{ho2020denoising,wavegrad}. Recent works have proposed improvements for diffusion model training, including importance sampling of the ELBO, better noise schedules \citep{nichol2021improved} and implicit diffusion models \citep{song2021implicit}. Several works have also drawn connections to score matching \citep{vincent2011connection,hyvarinen2004independent,song_2019}, leading to improved sampling algorithms in the continuous-time limit \citep{song2020_sde}.

While most works have considered continuous diffusion models, discrete diffusion-like models were described in \citep{sohl2015deep} and applied to text generation and image segmentation data in \citep{hoogeboom2021argmax}. Some works \citep{permutation,music} have dealt with discrete data by embedding it in continuous space and leveraging Gaussian diffusion, but have not applied this to text. \citet{princeton_chem} also considered generation of discrete structured objects using a diffusion-like Markov corruption process.

For text, denoising autoencoders have a long history both in representation learning \citep{Bengio_2013, bert} and more recently as generative models \citep{bert_speak}. These closely resemble our absorbing state diffusion variants for a particular schedule and transition matrix (see Section \ref{sec:bert_connections}), although our framing allows us to compute log-likelihoods and experiment with alternative transition matrices. Other works have considered non-autoregressive translation and speech transcription via insertion and deletion \citep{levenstein,ruis2020insertion}, masking \citep{mask_predict}, and iteratively-refined sequence alignments \citep{chan2020imputer,saharia2020latentalignments}.